% ID: SEL-D0032-OB-E01
% ESTADO: cotejada. Segunda pasada visual y borrador compilado revisados.
% Fuente: Selectividad/matematicas_aplicadas/titular_junio_EXAMEN 5-ANDALUCÍA.pdf
% SHA-256: 5f24b200fb6cf34f6d72016ec3dc35787b4074de9b7752677e5b0d959ed0dfcd
% Página: 2; opción B; ejercicio original 1.

\textbf{(2.5 puntos)} Un distribuidor de software informático tiene en su cartera de clientes tanto a empresas como a particulares. Ha de conseguir al menos 25 empresas como clientes y el número de clientes particulares deberá ser como mínimo el doble que el de empresas. Por razones de eficiencia del servicio postventa, tiene estipulado un límite global de 120 clientes anuales. Cada empresa le produce 386 \texteuro{} de beneficio, mientras que cada particular le produce 229 \texteuro{}. \textquestiondown{}Qué combinación de empresas y particulares le proporcionará el máximo beneficio? \textquestiondown{}A cuánto ascenderá ese beneficio?
